\chapter{Operational Reference}
\label{chap:operations}

\section{Source-Code Architecture}
The implementation separates orchestration, semantics, detectors, providers,
and reporting so that a provider or serializer change does not silently alter
parsing. The following table maps the principal packages to their
responsibilities.

\begin{longtable}{p{5cm}p{9.2cm}}
\caption{Principal implementation modules.}\\
\toprule Source area & Responsibility \\ \midrule\endhead
\code{cli.py}, \code{api/} & Command-line and programmatic entry points,
argument/configuration resolution, mode selection, exit status and hand-off to
the runtime. \\
\code{runtime.py} & Scan orchestration, provider isolation, result assembly,
health evaluation, validation and output dispatch. \\
\code{file_manifest.py} & Reproducible discovery, selection, exclusions,
population identity and file/LOC denominators. \\
\code{parser.py}, \code{document.py} & Tolerant YAML loading, source-coordinate
retention, document shape and artifact classification. \\
\code{jinja_semantics.py} & Jinja name extraction, lexical locals, guard and
filter recognition, and uncertainty reasons. \\
\code{task_semantics.py} & Task/module/action identification, registration,
facts, loops, conditions, environment and execution-order evidence. \\
\code{symbol_index.py} & Repository definitions/uses, static include/import
edges, reachability and ordered persistent-definition propagation. \\
\code{role_context.py}, \code{include_vars_semantics.py} & Role defaults/vars,
dependencies, inventory-like sources and included-variable provenance. \\
\code{detectors/} & Review-checklist, paper-derived, partial and currently
non-automatable-boundary detector families. \\
\code{external/ansible_lint.py} & Project-root discovery/collapse, subprocess
execution, configuration/suppression handling and native evidence conversion. \\
\code{external/gitleaks.py} & Secret scan execution, path/fingerprint
normalisation, exclusions, redaction and evidence conversion. \\
\code{profiles.py} & Versioned rule/provider populations and profile invariants. \\
\code{report/} & Catalogue, canonical occurrences, CSV/JSON/SARIF/text/Markdown
serialisation, comparison identities, schema loading and validation. \\
\code{diff_workflow.py} & Local/commit/PR snapshot selection and safe
introduced/fixed/unchanged/moved comparison. \\
\bottomrule
\end{longtable}

\section{End-to-End Pipeline}
The pipeline is deliberately staged so that every final number can be traced
back to a repository, revision, selected file, detector, and provider run.

\begin{enumerate}
  \item \textbf{Repository acquisition.} Resolve the repository identity and
  freeze the commit SHA. Research scans use the recorded manifest rather than
  silently moving to the latest default branch.
  \item \textbf{Discovery.} Apply ignore rules, classify candidate extensions,
  and record manifest and selected populations separately. The distinction
  explains why the frozen corpus contains 64,876 YAML files but only 27,709
  selected Ansible artifacts.
  \item \textbf{Parsing.} Load YAML tolerantly, retain exact coordinates, scan
  source text where formatting matters, and extract embedded Jinja expressions.
  A parse problem becomes a diagnostic tied to its file and stage.
  \item \textbf{Artifact classification.} Decide whether a document is a
  playbook, task list, handler list, vars/defaults, inventory, metadata, or an
  import-only artifact. Only compatible detectors are scheduled.
  \item \textbf{Local analysis.} Run syntax-, task-, play-, file-, and
  directory-level rules and collect definition/use/guard evidence.
  \item \textbf{Repository model.} Build role context, symbol indexes, static
  include/import edges, handler relations, variable provenance, and execution
  order. Re-run graph-aware predicates with repository evidence.
  \item \textbf{External providers.} Execute ansible-lint at collapsed project
  roots and GitLeaks at the repository boundary. Preserve native rule IDs,
  fingerprints, messages, tool versions, and health.
  \item \textbf{Normalisation.} Convert heterogeneous records into the v3
  finding shape, redact sensitive material, and attach source/analysis
  population identities.
  \item \textbf{Canonicalisation.} Group only sufficiently specific evidence,
  select a lead occurrence, merge provider evidence, and retain all raw records.
  \item \textbf{Validation and export.} Validate the envelope and critical
  nested types against the packaged schema, copy that schema into the output,
  then produce JSON, CSV, SARIF, text, and Markdown projections.
  \item \textbf{Comparison.} For base/head or version comparisons, match exact
  IDs first, then semantic identities that account for file paths. Report introduced,
  fixed, unchanged, and moved/uncertain partitions.
\end{enumerate}

\section{Provider Lifecycle and Failure Semantics}
\begin{figure}[H]
\centering
\resizebox{\textwidth}{!}{%
\begin{tikzpicture}[
 s/.style={draw=ansblue,rounded corners,minimum width=28mm,minimum height=9mm,
 align=center,fill=anslightblue,font=\small},
 d/.style={draw=ansorange,diamond,aspect=2,align=center,fill=anslightyellow,font=\small},
 a/.style={-{Latex[length=2mm]},thick},node distance=8mm and 13mm]
\node[s] (plan) {Provider planned};
\node[d,right=of plan] (available) {Dependency available?};
\node[s,right=of available] (run) {Execute isolated run};
\node[d,below=of run] (valid) {Valid complete output?};
\node[s,left=of valid] (partial) {Partial / failed diagnostic};
\node[s,right=of valid] (complete) {Completed evidence};
\node[s,below=of available] (missing) {Unavailable diagnostic};
\draw[a] (plan)--(available);
\draw[a] (available)--node[above,font=\scriptsize]{yes}(run);
\draw[a] (available)--node[left,font=\scriptsize]{no}(missing);
\draw[a] (run)--(valid);
\draw[a] (valid)--node[above,font=\scriptsize]{no}(partial);
\draw[a] (valid)--node[above,font=\scriptsize]{yes}(complete);
\end{tikzpicture}
}
\caption{External-provider state machine. No failure path produces a clean result.}
\end{figure}

The analysis-health envelope distinguishes \code{completed}, \code{partial},
\code{failed}, and \code{unavailable} outcomes. Diagnostics record stage,
category, status, failure nature, message, and optional file. A missing binary
is environment-dependent. A tool timeout or incompatible repository
configuration is a provider execution outcome. A YAML parser exception is an
analysis diagnostic. None of these becomes an ordinary source smell, because
doing so would blame the repository for the analysis environment.

\subsection{ansible-lint}
AnsibleGuard delegates established Ansible lint checks to ansible-lint and
does not rename them as internally invented detectors. Native IDs such as FQCN,
syntax, risky permissions, or command-use rules remain in provider evidence.
Project roots are discovered from configuration and repository structure.
Nested unconfigured roots covered by an ancestor collapse into that ancestor;
independently configured nested roots remain separate. This prevents duplicate
linting without erasing configuration boundaries. ansible-lint skip/warn lists
and inline \code{noqa} suppressions are respected.

\subsection{GitLeaks}
GitLeaks supplies pattern and entropy evidence for content and optional history.
AnsibleGuard normalises paths and coordinates, preserves the native rule and
fingerprint, redacts the secret value, and attaches contextual status. A
provider match proves that the configured GitLeaks predicate fired; it does not
prove that an example token is active. This distinction explains why all 1,049
canonical GitLeaks occurrences in the frozen scan retain
\code{secret_context=uncertain_secret}.

\section{Exact v0.0.6 Profiles}
Profiles define a reproducible population of rules and providers. They are not
mere display filters: some require external tools, a manifest, or review-range
input.
These 11 named profiles were defined in AnsibleGuard for distinct users and
decisions; their rule subsets and provider choices are part of the system
design, rather than output groupings selected after the results were known.

\begin{longtable}{p{3.4cm}r p{3.4cm}p{6cm}}
\toprule Profile & Rules & Providers & Purpose and boundary \\ \midrule\endhead
\code{technical-full} & 76 & internal, ansible-lint, GitLeaks & All source,
external-tool, and Git-review findings. \\
\code{developer-fast} & 24 & internal & Fast security, correctness and
reliability feedback without external dependencies. \\
\code{security} & 9 & all three & Eight security concepts plus the selected
security-relevant ansible-lint subset and GitLeaks. \\
\code{lint} & 1 & ansible-lint & Native ansible-lint findings only, carried by
the lint envelope and native evidence. \\
\code{quality} & 54 & internal & Correctness, dependencies, documentation,
maintainability, reliability, structure, style and testing. \\
\code{pr} & 69 & all three & Base/head source and provider snapshots; the seven
review-metadata rules run in their separate range phase. \\
\code{snapshot-source} & 67 & internal & Frozen source state excluding review
range and external providers; requires a manifest. \\
\code{research-core} & 8 & internal & Conservative paper-derived allowlist;
provisional pending manual calibration. \\
\code{extended-ansible} & 59 & internal & Disjoint Ansible-specific secondary
stratum outside the research core; requires a manifest. \\
\code{review-range} & 7 & internal & Commit/base-to-head review-process rules
only; no source snapshot analysis. \\
\code{incremental-file-local} & 6 & internal & Explicitly partial changed-file
feedback; never equivalent to a full repository comparison. \\
\bottomrule
\end{longtable}

\begin{figure}[htbp]
\centering
\begin{tikzpicture}[x=1cm,y=1cm]
  \fill[anspaper,rounded corners=4pt] (0,0) rectangle (14.4,3.5);
  \node[anchor=west,text=ansnavy,font=\bfseries\large] at (.35,3.06)
    {Profiles are designed around decisions};
  \node[draw=ansblue,fill=white,rounded corners=4pt,
    text width=3.8cm,minimum height=2.05cm,align=left,anchor=north west,
    inner sep=7pt] at (.35,2.65)
    {\textcolor{ansblue}{\bfseries Authors}\\[3pt]
    \small Local feedback\\[3pt]
    \footnotesize\code{developer-fast}\\\code{quality}};
  \node[draw=ansteal,fill=white,rounded corners=4pt,
    text width=3.8cm,minimum height=2.05cm,align=left,anchor=north west,
    inner sep=7pt] at (5.0,2.65)
    {\textcolor{ansteal}{\bfseries Reviewers}\\[3pt]
    \small PR and security evidence\\[3pt]
    \footnotesize\code{pr} / \code{security}\\\code{review-range}};
  \node[draw=ansgold,fill=white,rounded corners=4pt,
    text width=3.8cm,minimum height=2.05cm,align=left,anchor=north west,
    inner sep=7pt] at (9.65,2.65)
    {\textcolor{ansnavy}{\bfseries Researchers}\\[3pt]
    \small Frozen populations\\[3pt]
    \footnotesize\code{technical-full}\\\code{research-core}};
\end{tikzpicture}
\caption{Representative intended users and profiles. The complete eleven-profile
definition, including manifest and provider requirements, is given above.}
\label{fig:profile_audiences}
\end{figure}

The eight \code{research-core} concepts are \code{complex-expression},
\code{incomplete-task}, \code{long-statement}, \code{unstructured-module},
\code{empty-abstraction}, \code{excessive-nesting},
\code{missing-abstraction}, and \code{broken-role-dependency}. The six
incremental file-local rules are \code{long-task-name}, \code{large-task},
\code{plaintext-sensitive-variable}, \code{missing-validate-certs},
\code{loop-legacy-syntax}, and \code{variable-naming-convention}.

\subsection{Rule Layers Used in the 149-Repository Study}
The frozen study uses \code{technical-full}, the only profile intended to
inventory all 76 catalogue concepts together with both external providers.
``All rules'' does not mean that every detector executes on every file. The
artifact classifier and rule metadata restrict each predicate to compatible
contexts, as summarised below; the complete rule-by-rule predicates appear in
Appendix~\ref{chap:rules}.

\begin{longtable}{@{}p{3.0cm}p{3.3cm}p{3.8cm}p{3.7cm}@{}}
\caption{Analysis-rule layers applied to each frozen repository.}
\label{tab:study_rule_layers}\\
\toprule Layer & Owner & Applicable evidence & Study role \\
\midrule\endfirsthead
\toprule Layer & Owner & Applicable evidence & Study role \\
\midrule\endhead
Source syntax and structure & AnsibleGuard & Selected playbooks, tasks,
handlers, variable/default, inventory, metadata and role/collection artifacts
according to rule compatibility & Detect malformed, oversized, incomplete, or
misplaced structures while avoiding unrelated YAML. \\
Repository semantics & AnsibleGuard & Cross-file definitions, includes/imports,
roles, handlers, dependencies, documentation and test layout & Detect concerns
that a file-local linter cannot decide safely. \\
Variable/runtime evidence & AnsibleGuard & Jinja uses, static definitions,
persistent runtime producers, guards, lexical locals, ordering and named
dynamic sources & Assign \code{CONFIRMED}, \code{UNCERTAIN},
\code{UNRESOLVED}, or \code{ORDERED_CONFLICT} with traceable evidence. \\
Established Ansible lint & ansible-lint & Each resolved project root under its
own configuration and suppression policy & Delegate FQCN, naming, YAML and
module-practice rules without pretending they are native inventions. \\
Candidate secrets & GitLeaks & Selected repository content under provider
configuration, with secret value redacted & Supply pattern/entropy evidence and
native rule identity for security canonicalisation. \\
Review-range metadata & AnsibleGuard & Git history/base--head or PR metadata
when the required range exists & Seven range concepts are in the 76-concept
profile, but require suitable Git review context; absence of a range is not a
source finding. \\
Canonical overlap & AnsibleGuard evidence layer & Normalised native,
ansible-lint and GitLeaks records & Merge one review concern without deleting
provider provenance or cross-file identity safeguards. \\
\bottomrule
\end{longtable}

Thus, the 149-repository result is the composition of compatible native rules,
delegated linter evidence, candidate-secret evidence, and any available
range-dependent evidence under one validated contract. Provider failure or a
missing range changes analysis health/availability; it does not manufacture an
ordinary source smell.

\section{Output Directory and File Responsibilities}
\begin{longtable}{p{4.2cm}p{10cm}}
\toprule Path family & Contents and use \\ \midrule\endhead
\code{summary/} & Per-repository and all-repository CSV/JSON aggregations,
including population denominators, provider health, diagnostics, and grouped
counts. \\
\code{raw/} & Provider-specific records before cross-provider canonicalisation;
used to reproduce transformations and debug adapters. \\
\code{canonical/} & Canonical occurrences, lead detector, merged evidence,
semantic identity, status and severity. This is the principal analysis unit. \\
\code{metadata/} & Repository URL/slug, commit SHA, tool and dependency
versions, configuration/profile, manifest identity, timing and diagnostics. \\
\code{files/} & Optional per-file projections, maintaining repository-relative
paths and source regions. \\
\code{diff/} & Comparison partitions and matching evidence for base/head,
commit, PR, or release comparisons. \\
\code{schemas/output-v3.json} & Exact JSON Schema copied from the installed
package and used to validate this output. \\
\bottomrule
\end{longtable}

\section{Versioned Envelope}
Every versioned document identifies \code{schema_version}, \code{kind}, and a
\code{contract}. Depending on kind, it contains \code{items}, aggregate
\code{counts}, \code{analysis_health}, \code{artifacts}, \code{policy}, or a
\code{comparison}. The shared validator rejects a malformed versioned envelope.
The PR comparison script accepts an explicit legacy list only when the input is
actually a list; invalid v3 data cannot fall through as legacy.

\subsection{Finding fields}
\begin{longtable}{p{4cm}p{10.2cm}}
\toprule Field & Meaning \\ \midrule\endhead
\code{finding_id} & Stable provider-level record identity. \\
\code{provider} & Producing component: internal, ansible-lint, or GitLeaks. \\
\code{rule_id} & Normalised AnsibleGuard concept/envelope. \\
\code{concept_id} & Cross-provider conceptual identity used for overlap. \\
\code{file}, \code{line} & Repository-relative source location and positive line. \\
\code{severity} & Policy level: error, warning, or note. \\
\code{message} & Human-readable explanation after redaction. \\
\code{source_population} & Identity of the source/manifest population. \\
\code{analysis_population} & Identity of the rule/provider population. \\
\bottomrule
\end{longtable}

\subsection{Canonical occurrence fields}
\begin{longtable}{p{4cm}p{10.2cm}}
\toprule Field & Meaning \\ \midrule\endhead
\code{occurrence_id} & Stable canonical record identity. \\
\code{concept_id} & Normalised concept shared by overlapping detectors. \\
\code{file}, \code{line} & Lead source location; cross-file matching requires
specific file-aware evidence. \\
\code{detectors} & Set of detectors contributing evidence. \\
\code{semantic_key} & Rule-aware identity for controlled comparison. \\
\code{content_semantic_key} & Content-based identity used only under safe
matching constraints. \\
\code{semantic_subject} & Variable, task, handler, role, secret class, or other
subject that the occurrence describes. \\
\code{task_fingerprint} & Stable task identity where task structure exists. \\
\code{provider_evidence_fingerprint} & Digest of merged native evidence. \\
\code{scan_mode} & Full, snapshot, PR, commit, or explicitly partial mode. \\
\code{source_population}, \code{analysis_population} & Denominators needed for
valid comparisons. \\
\bottomrule
\end{longtable}

\subsection{Provider-evidence and diagnostic fields}
Provider evidence contains \code{provider}, \code{rule_id},
\code{native_rule_id}, \code{finding_id}, \code{native_fingerprint}, and
\code{message}. Diagnostics contain \code{stage}, \code{category},
\code{status}, \code{failure_nature}, \code{message}, and optional
\code{file}. Analysis health contains an aggregate \code{status} and boolean
\code{complete}. These fields make a partial provider distinguishable from a
complete zero-finding run.

\section{Meaning of Summary CSV Columns}
The exact wide summary may add one column per category, status, severity,
provider, or concept, but all columns belong to one of the following groups.

\begin{longtable}{p{4.4cm}p{9.8cm}}
\toprule Column group & Interpretation \\ \midrule\endhead
Repository columns & Stable slug/URL, recorded SHA and completion state; define
which repository revision was measured. \\
Manifest files/LOC & All YAML candidates in the frozen manifest; broad corpus
denominator. \\
Selected files/LOC & Files classified as in-scope Ansible artifacts and their
line count; actual analysis denominator. \\
Raw findings & Sum emitted by providers before canonical merge. \\
Canonical occurrences & Deduplicated evidence units; not the sum of category
columns if a projection intentionally excludes a zero/unavailable group. \\
Provider raw columns & Native contribution from AnsibleGuard, ansible-lint, and
GitLeaks before merge. \\
Provider lead columns & Provider selected as lead for a canonical occurrence;
not the same as total evidence contributed by that provider. \\
Status columns & CONFIRMED, UNCERTAIN, EXPOSED, UNRESOLVED, and
\code{ORDERED_CONFLICT} detector-state counts. \\
Severity columns & Error, warning, and note policy-impact counts. \\
Category columns & Canonical occurrences grouped by taxonomy category. \\
Concept/rule columns & Occurrence counts for each normalised concept or native
rule projection. \\
Provider health columns & Completed, partial, failed, unavailable and diagnostic
counts; required before interpreting zeroes. \\
Timing/version columns & Runtime, package/provider versions, profile and schema
identity for reproduction. \\
\bottomrule
\end{longtable}

\section{Full, Local, Commit, and Pull-Request Operation}
\textbf{Full repository mode} scans the selected manifest, constructs the
repository graph, invokes configured providers, and writes a complete output
tree. \textbf{Local mode} can provide quick changed-file feedback but must say
when repository-global and external checks are deferred. \textbf{Commit mode}
resolves two revisions, analyses both compatible populations, then compares
canonical evidence. \textbf{PR mode} uses the base/head revisions and adds the
review-range rules for commit message, dependency trailers, release notes,
copyright, and test-update obligations.

The safe comparison hierarchy is: exact occurrence identity; native provider
fingerprint; file-aware semantic/task identity; then explicitly uncertain
movement. Generic line/column findings cannot cross file boundaries and become
semantic renames. This prevents a common error in which unrelated findings on
the same line number in two files are reported as one moved occurrence.

\section{Version Evolution and Evidence}
This section provides historical traceability. The current v0.0.7 study is in
Chapter~\ref{chap:v007}; earlier releases appear where they explain a
requirement, regression test, or controlled change.
\begin{longtable}{p{2.2cm}p{5.5cm}p{6.6cm}}
\toprule Version & Problem or objective & Implemented response and evidence \\
\midrule\endhead
v0.0.4 & Establish a frozen full-study baseline and identify provider and
semantic weaknesses. & Produced the baseline population used for later paired
comparisons; exposed the need to separate provider health and canonical data. \\
v0.0.5 & Prevent provider crashes from aborting repositories and stop partial
ansible-lint runs from appearing clean. & Added provider isolation, explicit
partial diagnostics, canonical occurrence output, GitLeaks integration, and
initial semantic uncertainty. A 15-repository diagnostic replay completed all
repositories but still produced 26,552 unguarded-variable occurrences, so it
was not accepted as a semantic endpoint. \\
v0.0.6 & Explain remaining variable findings and harden reproducibility. & Added
static include/import propagation without disconnected-file leakage, lexical
Jinja loop scope, safer occurrence matching, shared v3 validation, packaged
schema copies, stronger field/type checks, ansible-lint root-collapse tests,
and targeted rule fixes. The frozen 149-repository scan completed and its
mechanism distributions are reported in Chapter~\ref{chap:results}. \\
\bottomrule
\end{longtable}

Changes in totals have more than one cause: restored identity can increase
canonical occurrences; provider suppression can lower raw findings; new
uncertainty reasons can move cases out of non-uncertain states; stronger
canonicalisation can merge duplicates; and rule fixes can remove demonstrated
false-positive mechanisms. Therefore, no version delta is interpreted without
same-population denominators and mechanism evidence.

\section{Reproduction and Validation Checklist}
For a defensible rerun, record the branch and AnsibleGuard commit, pull both
code and frozen manifests, create an isolated environment, install the package
and provider dependencies, verify provider versions, run the selected profile,
and monitor per-repository logs. After execution, check 149/149 repository
states, provider-health totals, diagnostics, selected files/LOC, schema
validation, and summary reconciliation before interpreting any rule count.

Release validation includes focused regression tests, the full pytest suite,
checks for stale documentation, Python byte-code compilation, whitespace checks,
optional remote equivalence tests, and available benchmark gates.

Research validation adds comparison on a frozen population, mechanism sampling,
and independent manual labels before any precision claim. It also requires
adjudication, confidence intervals, and a seeded-defect corpus for recall.
